% GENERATED FILE. Edit canonical lessons, then run npm run generate:books.
% canonical-source-hash: fnv1a64:1bf9e3465aa83ba7
% canonical-lessons: FR-C320-trinquer, FR-C320-toast, FR-C320-souhaiter, FR-C320-bougie, FR-C320-voeu

\chapter{À votre santé !}
\label{ch:fr-a-votre-sante}
\hlchaptermodality{320}

\begin{chapteropening}
\textbf{By the end of this chapter:} \emph{I can raise a glass and make wishes: clinking glasses, a toast, wishing, candles and a wish.}

Complete the last lesson of chapter 320: I can raise a glass and make wishes: clinking glasses, a toast, wishing, candles and a wish.
\end{chapteropening}

\section[trinquer]{trinquer — to clink glasses, to drink to}

\label{lesson:FR-C320-trinquer}



\begin{quote}
Before the new one: say the French for the log, then the French for the king cake.
\end{quote}

\subsection*{You'll want to know: trinquer}
\textbf{trinquer} — \textquotedblleft{}to clink glasses, to drink to\textquotedblright{}.

\begin{quote}
Nous avons trinqué aux mariés.
\end{quote}

\begin{quote}
\emph{We clinked glasses to the bride and groom.}
\end{quote}

Informal: \emph{c'est moi qui ai trinqué} — I paid the price.

\begin{cousinweb}
\emph{trinquer} was borrowed from German \emph{trinken} \textquotedblleft{}to drink\textquotedblright{} (EtymDB).
\end{cousinweb}

\begin{grammarlens}[title={What you've built}]
One more word for telling, explaining and describing at B1.
\end{grammarlens}

\subsection*{Your turn}
\begin{itemize}
\raggedright
  \item \emph{Say it:} \emph{trinquer}
  \item \emph{Say it:} \emph{trinquer}, once more
  \item \emph{Say it:} a sentence of your own with \emph{trinquer}
  \item \emph{Recall:} say the French for the log, then the French for the king cake, then say \emph{trinquer} again
\end{itemize}

\begin{tcolorbox}[breakable,colback=teal!4,colframe=teal!35!black,title={Before you move on}]
What does trinquer mean? (\textquotedblleft{}To clink glasses, to drink to\textquotedblright{}.) Say it once more.
\end{tcolorbox}
\section[le toast]{le toast — the toast (with a drink); toasted bread}

\label{lesson:FR-C320-toast}



\begin{quote}
Before the new one: say the French for the king cake, then the French for to clink glasses.
\end{quote}

\subsection*{You'll want to know: le toast}
\textbf{le toast} — \textquotedblleft{}the toast (with a drink); toasted bread\textquotedblright{}.

\begin{quote}
Le père de la mariée a porté un toast.
\end{quote}

\begin{quote}
\emph{The bride's father proposed a toast.}
\end{quote}

\emph{Porter un toast à quelqu'un} — to drink to someone. Also a slice of toasted bread.

\begin{cousinweb}
\emph{toast} is usually traced to English \emph{toast} (EtymDB lists no parent), which comes from Old French \emph{toster} \textquotedblleft{}to grill\textquotedblright{}, built on Latin \emph{tostus} \textquotedblleft{}roasted\textquotedblright{}, from \emph{torrēre} (Lewis \& Short).
\end{cousinweb}

\begin{grammarlens}[title={What you've built}]
One more word for telling, explaining and describing at B1.
\end{grammarlens}

\subsection*{Your turn}
\begin{itemize}
\raggedright
  \item \emph{Say it:} \emph{le toast}
  \item \emph{Say it:} \emph{le toast}, once more
  \item \emph{Say it:} a sentence of your own with \emph{le toast}
  \item \emph{Recall:} say the French for the king cake, then the French for to clink glasses, then say \emph{le toast} again
\end{itemize}

\begin{tcolorbox}[breakable,colback=teal!4,colframe=teal!35!black,title={Before you move on}]
What does le toast mean? (\textquotedblleft{}The toast (with a drink); toasted bread\textquotedblright{}.) Say it once more.
\end{tcolorbox}
\section[souhaiter]{souhaiter — to wish}

\label{lesson:FR-C320-souhaiter}



\begin{quote}
Before the new one: say the French for to clink glasses, then the French for the toast.
\end{quote}

\subsection*{You'll want to know: souhaiter}
\textbf{souhaiter} — \textquotedblleft{}to wish\textquotedblright{}.

\begin{quote}
Je te souhaite un joyeux anniversaire.
\end{quote}

\begin{quote}
\emph{I wish you a happy birthday.}
\end{quote}

\emph{Souhaiter que} + subjunctive: \emph{je souhaite qu'il vienne}.

\begin{cousinweb}
\emph{souhaiter} comes from Old French \emph{sohaidier} (EtymDB), built with \emph{sub-} on a Germanic verb meaning \textquotedblleft{}to promise, to command\textquotedblright{}.
\end{cousinweb}

\begin{grammarlens}[title={What you've built}]
One more word for telling, explaining and describing at B1.
\end{grammarlens}

\subsection*{Your turn}
\begin{itemize}
\raggedright
  \item \emph{Say it:} \emph{souhaiter}
  \item \emph{Say it:} \emph{souhaiter}, once more
  \item \emph{Say it:} a sentence of your own with \emph{souhaiter}
  \item \emph{Recall:} say the French for to clink glasses, then the French for the toast, then say \emph{souhaiter} again
\end{itemize}

\begin{tcolorbox}[breakable,colback=teal!4,colframe=teal!35!black,title={Before you move on}]
What does souhaiter mean? (\textquotedblleft{}To wish\textquotedblright{}.) Say it once more.
\end{tcolorbox}
\section[la bougie]{la bougie — the candle; the spark plug}

\label{lesson:FR-C320-bougie}



\begin{quote}
Before the new one: say the French for the toast, then the French for to wish.
\end{quote}

\subsection*{You'll want to know: la bougie}
\textbf{la bougie} — \textquotedblleft{}the candle; the spark plug\textquotedblright{}.

\begin{quote}
Elle a soufflé les trente bougies de son gâteau.
\end{quote}

\begin{quote}
\emph{She blew out the thirty candles on her cake.}
\end{quote}

In an engine, \emph{une bougie} is a spark plug.

\begin{cousinweb}
\emph{bougie} comes from the name of a North African port town that sold candle wax to Europe in the Middle Ages (EtymDB).
\end{cousinweb}

\begin{grammarlens}[title={What you've built}]
One more word for telling, explaining and describing at B1.
\end{grammarlens}

\subsection*{Your turn}
\begin{itemize}
\raggedright
  \item \emph{Say it:} \emph{la bougie}
  \item \emph{Say it:} \emph{la bougie}, once more
  \item \emph{Say it:} a sentence of your own with \emph{la bougie}
  \item \emph{Recall:} say the French for the toast, then the French for to wish, then say \emph{la bougie} again
\end{itemize}

\begin{tcolorbox}[breakable,colback=teal!4,colframe=teal!35!black,title={Before you move on}]
What does la bougie mean? (\textquotedblleft{}The candle; the spark plug\textquotedblright{}.) Say it once more.
\end{tcolorbox}
\section[le vœu]{le vœu — the wish; the vow}

\label{lesson:FR-C320-voeu}



\begin{quote}
Before the new one: say the French for to wish, then the French for the candle.
\end{quote}

\subsection*{You'll want to know: le vœu}
\textbf{le vœu} — \textquotedblleft{}the wish; the vow\textquotedblright{}.

\begin{quote}
Elle a fait un vœu avant de souffler les bougies.
\end{quote}

\begin{quote}
\emph{She made a wish before blowing out the candles.}
\end{quote}

\emph{Meilleurs vœux !} — best wishes for the new year.

\begin{cousinweb}
\emph{vœu} comes from Old French \emph{vou} (EtymDB), inherited from Latin \emph{vōtum} \textquotedblleft{}a vow, a wish\textquotedblright{}, from \emph{vovēre} \textquotedblleft{}to vow\textquotedblright{} (Lewis \& Short).
\end{cousinweb}

\begin{grammarlens}[title={What you've built}]
One more word for telling, explaining and describing at B1.
\end{grammarlens}

\subsection*{Your turn}
\begin{itemize}
\raggedright
  \item \emph{Say it:} \emph{le vœu}
  \item \emph{Say it:} \emph{le vœu}, once more
  \item \emph{Say it:} a sentence of your own with \emph{le vœu}
  \item \emph{Recall:} say the French for to wish, then the French for the candle, then say \emph{le vœu} again
\end{itemize}

\begin{tcolorbox}[breakable,colback=teal!4,colframe=teal!35!black,title={Before you move on}]
What does le vœu mean? (\textquotedblleft{}The wish; the vow\textquotedblright{}.) Say it once more.
\end{tcolorbox}
